\begin{proof}[Proof of \cref{obj:lem:uniform-affine-profile-bound}]
Fix an admissible \(v\), an integer \(n\ge2\), and \(P\in\mathcal M_v\). Write
\[
\mu=P^{\otimes n},
\qquad
M=J_{\mathrm c}.
\]
All expectations below are under \(\mu\).

\begin{enumerate}
\item
If \(n<4\), then \(n\in\{2,3\}\). The construction in \cref{obj:def:explicit-score-ledger} gives, for every sample and every \(c\in\mathbb R\),
\[
Z_1(c)=Z_2(c)=0,
\qquad
\theta_P(c)=0,
\qquad
W(c)=0,
\qquad
\Lambda_{n,v}=0.
\]
Consequently, the asserted profile event is the entire sample space, with probability one, and the equality assertion holds for every sample. Henceforth suppose \(n\ge4\).
% lean: CausalSmith.Stat.FinitepHomogeneityDensegamma.profileBound_small_sample

\item
Define the profile discrepancy and its simultaneous event by
\[
\mathsf F_{\mathrm{prof}}(c)
=
\bigl|W(c)-\|\theta_P(c)\|_2^2\bigr|
-
128\left(
\sqrt{\Lambda_{n,v}}\,\|\theta_P(c)\|_2
+\sqrt M\,\Lambda_{n,v}
\right),
\]
\[
\mathcal G_{\mathrm{prof}}
=
\left\{
\mathsf F_{\mathrm{prof}}(c)\le0
\text{ for every }c\in[-1/2,1/2]
\right\}.
\]
For fixed \(c\), the discrepancy is a measurable function of the sample. For each fixed sample, it is continuous in \(c\), because the scores and their population means are affine. Choose a countable dense set with
\[
\mathcal C_{\mathrm{dense}}\subseteq[-1/2,1/2],
\qquad
\overline{\mathcal C_{\mathrm{dense}}}=[-1/2,1/2].
\]
Continuity gives
\[
\mathcal G_{\mathrm{prof}}
=
\bigcap_{c\in\mathcal C_{\mathrm{dense}}}
\{\mathsf F_{\mathrm{prof}}(c)\le0\}.
\]
Indeed, for each sample, the set where the discrepancy is nonpositive is closed in the profiling interval. Thus \(\mathcal G_{\mathrm{prof}}\) is measurable.

We also record the positivity needed for the tail estimates and the zero-ledger assertion. The public formulas in \cref{obj:def:explicit-score-ledger} give
\[
s=s_n=\lfloor n/2\rfloor\ge2,
\qquad
M\ge1,
\qquad
T_0\ge1,
\qquad
V_0=32T_0^{2-p}>0.
\]
In the single-resolution branch, \(L_1=16\sqrt{V_0}\). In the multiresolution branch, the positive ranks and cutoffs give
\[
a_j\ge0,\qquad d_j\ge0,\qquad V_j>0
\quad(1\le j\le L),
\]
and hence
\[
\sum_{j=1}^{L}\bigl(d_j+a_j\sqrt{V_j}\bigr)\ge0,
\qquad
L_1
=
16\sqrt{V_0}
+8\sum_{j=1}^{L}\bigl(d_j+a_j\sqrt{V_j}\bigr)
\ge16\sqrt{V_0}>0.
\]
These inequalities include the empty-sum case \(L=0\). In either branch,
\[
\Lambda_{n,v}
=
\frac{L_1^2}{s}+\frac{2L_2^2}{s(s-1)}
\ge\frac{L_1^2}{s}>0.
\]
Together with the small-sample branch, this establishes, for admissible \(v\) and \(n\ge2\),
\[
\Lambda_{n,v}=0
\quad\Longleftrightarrow\quad
n<4.
\]
The zero-ledger assertion therefore follows from the pointwise equalities in the first case.
% lean: CausalSmith.Stat.FinitepHomogeneityDensegamma.measurableSet_profileEvent
% lean: CausalSmith.Stat.FinitepHomogeneityDensegamma.ledgerCov_pos_of_large_sample
% lean: CausalSmith.Stat.FinitepHomogeneityDensegamma.ledger_profile_zero_covariance

\item
For \(b\in\{1,2\}\), define the affine coefficients by
\[
Z_{b,0}=Z_b(0),
\qquad
Z_{b,A}=Z_b(0)-Z_b(1),
\qquad
Z_b(c)=Z_{b,0}-cZ_{b,A}
\quad(c\in\mathbb R).
\]
The common block-mean identity supplied by \cref{obj:lem:original-score-mean-ledger} is
\[
\mathbb E_\mu[Z_b(c)]=\theta_P(c)
\quad(b\in\{1,2\},\ c\in\mathbb R).
\]
Define the population coefficients and centered coefficient errors by
\[
\overline Z_0=\theta_P(0),
\qquad
\overline Z_A=\theta_P(0)-\theta_P(1),
\]
\[
\mathsf E_{b,\iota}=Z_{b,\iota}-\overline Z_\iota
\quad(b\in\{1,2\},\ \iota\in\{0,A\}),
\qquad
\mathsf E_b(c)=\mathsf E_{b,0}-c\mathsf E_{b,A}.
\]
Linearity of expectation gives
\[
\mathbb E_\mu[Z_{b,\iota}]=\overline Z_\iota,
\qquad
\mathbb E_\mu[\mathsf E_{b,\iota}]=0,
\]
\[
\theta_P(c)=\overline Z_0-c\overline Z_A,
\qquad
Z_b(c)=\theta_P(c)+\mathsf E_b(c).
\]
Introduce the fixed population subspace
\[
\mathcal S_P=\operatorname{span}\{\overline Z_0,\overline Z_A\}
\subseteq\mathbb R^M.
\]
Its orthogonal projection \(\operatorname{pr}_{\mathcal S_P}:\mathbb R^M\to\mathcal S_P\) is characterized by
\[
\operatorname{pr}_{\mathcal S_P}\mathsf v_{\mathrm{dir}}\in\mathcal S_P,
\qquad
\mathsf v_{\mathrm{dir}}-\operatorname{pr}_{\mathcal S_P}\mathsf v_{\mathrm{dir}}\in\mathcal S_P^\perp
\quad(\mathsf v_{\mathrm{dir}}\in\mathbb R^M).
\]

For \(b\in\{1,2\}\) and \(\iota,\jmath\in\{0,A\}\), define the bad events
\[
\mathcal P_{b,\iota}
=
\left\{
\|\operatorname{pr}_{\mathcal S_P}\mathsf E_{b,\iota}\|_2
>
40\sqrt{\Lambda_{n,v}}
\right\},
\]
\[
\mathcal Q_{\iota,\jmath}
=
\left\{
|\langle\mathsf E_{1,\iota},\mathsf E_{2,\jmath}\rangle|
>
40\sqrt M\,\Lambda_{n,v}
\right\},
\]
\[
\mathcal F_{\mathrm{bad}}
=
\left(
\bigcup_{b\in\{1,2\}}\ \bigcup_{\iota\in\{0,A\}}
\mathcal P_{b,\iota}
\right)
\cup
\left(
\bigcup_{\iota\in\{0,A\}}\ \bigcup_{\jmath\in\{0,A\}}
\mathcal Q_{\iota,\jmath}
\right).
\]
The coefficients are measurable, so continuity of projection, norm, and inner product makes all these events measurable.
% lean: CausalSmith.Stat.FinitepHomogeneityDensegamma.ledgerScore_eq_theta_add_error
% lean: CausalSmith.Stat.FinitepHomogeneityDensegamma.uniform_affine_profile_bound

\item
The square integrability and variance bounds in \cref{obj:lem:original-score-covariance-ledger}, together with the coefficient means just computed, imply
\[
\mathbb E_\mu\!\left[
\langle \mathsf v_{\mathrm{dir}},\mathsf E_{b,\iota}\rangle^2
\right]
=
\operatorname{Var}_\mu\!\left(\langle \mathsf v_{\mathrm{dir}},Z_{b,\iota}\rangle\right)
\le
\Lambda_{n,v}\|\mathsf v_{\mathrm{dir}}\|_2^2
\quad(\mathsf v_{\mathrm{dir}}\in\mathbb R^M).
\]
Because \(\mathcal S_P\) is spanned by two vectors,
\[
\dim\mathcal S_P\le2.
\]
Choose an orthonormal basis
\[
\mathcal B_P
=
\{\mathsf v^{\mathrm{sp}}_{\mathrm{dir},1},\ldots,\mathsf v^{\mathrm{sp}}_{\mathrm{dir},\dim\mathcal S_P}\}
\subseteq\mathcal S_P.
\]
Parseval's identity and the defining property of orthogonal projection give
\[
\begin{aligned}
\mathbb E_\mu
\|\operatorname{pr}_{\mathcal S_P}\mathsf E_{b,\iota}\|_2^2
&=
\sum_{\mathsf v_{\mathrm{dir}}\in\mathcal B_P}
\mathbb E_\mu
\langle \mathsf v_{\mathrm{dir}},\operatorname{pr}_{\mathcal S_P}\mathsf E_{b,\iota}\rangle^2\\
&=
\sum_{\mathsf v_{\mathrm{dir}}\in\mathcal B_P}
\mathbb E_\mu\langle \mathsf v_{\mathrm{dir}},\mathsf E_{b,\iota}\rangle^2\\
&\le
(\dim\mathcal S_P)\Lambda_{n,v}
\le2\Lambda_{n,v}.
\end{aligned}
\]
If \(\mathcal S_P=\{0\}\), the basis is empty and these sums equal zero. Markov's inequality applied to the squared projected norm now yields
\[
\mu(\mathcal P_{b,\iota})
\le
\frac{
\mathbb E_\mu
\|\operatorname{pr}_{\mathcal S_P}\mathsf E_{b,\iota}\|_2^2
}{
1600\Lambda_{n,v}
}
\le\frac1{800}.
\]
There are two blocks and two coefficients, so the finite union bound gives
\[
\mu\left(
\bigcup_{b\in\{1,2\}}\ \bigcup_{\iota\in\{0,A\}}
\mathcal P_{b,\iota}
\right)
\le
\sum_{b\in\{1,2\}}\ \sum_{\iota\in\{0,A\}}
\mu(\mathcal P_{b,\iota})
\le\frac4{800}.
\]
% lean: CausalSmith.Stat.FinitepHomogeneityDensegamma.ledgerCoeffError_inner_second_moment_le
% lean: CausalSmith.Stat.FinitepHomogeneityDensegamma.ledgerCoeffError_projection_tail

\item
The coefficients in the two blocks depend on disjoint sets of independent records, as specified in \cref{obj:def:blocks,obj:def:explicit-score-ledger}. Subtracting deterministic means preserves their independence. Define their probability laws on \(\mathbb R^M\) by
\[
\nu_{b,\iota}
=
\mu\circ\mathsf E_{b,\iota}^{-1}.
\]
For every \(\iota,\jmath\in\{0,A\}\), their joint law is therefore
\[
\mu\circ
(\mathsf E_{1,\iota},\mathsf E_{2,\jmath})^{-1}
=
\nu_{1,\iota}\otimes\nu_{2,\jmath}.
\]
Choose an orthonormal basis of the ambient space,
\[
\mathcal B_{\mathrm{amb}}
=
\{\mathsf v^{\mathrm{amb}}_{\mathrm{dir},1},\ldots,\mathsf v^{\mathrm{amb}}_{\mathrm{dir},M}\}
\subseteq\mathbb R^M.
\]
The directional second-moment bound from the preceding step gives
\[
\mathbb E_\mu\|\mathsf E_{1,\iota}\|_2^2
=
\sum_{\mathsf v_{\mathrm{dir}}\in\mathcal B_{\mathrm{amb}}}
\mathbb E_\mu\langle \mathsf v_{\mathrm{dir}},\mathsf E_{1,\iota}\rangle^2
\le M\Lambda_{n,v}.
\]
Square integrability and
\[
|\langle \mathsf v_{\mathrm{dir}},z\rangle|^2\le\|\mathsf v_{\mathrm{dir}}\|_2^2\|z\|_2^2
\quad(\mathsf v_{\mathrm{dir}},z\in\mathbb R^M)
\]
justify iterated integration under the product law. Applying the directional second-moment bound to the second coefficient, for each fixed \(\mathsf v_{\mathrm{dir}}\), gives
\[
\begin{aligned}
\mathbb E_\mu
\langle\mathsf E_{1,\iota},\mathsf E_{2,\jmath}\rangle^2
&=
\int_{\mathbb R^M}\int_{\mathbb R^M}
\langle \mathsf v_{\mathrm{dir}},z\rangle^2\,
d\nu_{2,\jmath}(z)\,d\nu_{1,\iota}(\mathsf v_{\mathrm{dir}})\\
&\le
\int_{\mathbb R^M}
\Lambda_{n,v}\|\mathsf v_{\mathrm{dir}}\|_2^2\,d\nu_{1,\iota}(\mathsf v_{\mathrm{dir}})\\
&=
\Lambda_{n,v}\,
\mathbb E_\mu\|\mathsf E_{1,\iota}\|_2^2\\
&\le M\Lambda_{n,v}^2.
\end{aligned}
\]
Since \(M\ge1\) and \(\Lambda_{n,v}>0\), Markov's inequality gives
\[
\mu(\mathcal Q_{\iota,\jmath})
\le
\frac{
\mathbb E_\mu
\langle\mathsf E_{1,\iota},\mathsf E_{2,\jmath}\rangle^2
}{
1600M\Lambda_{n,v}^2
}
\le\frac1{1600}.
\]
There are four coefficient pairs, whence
\[
\mu\left(
\bigcup_{\iota\in\{0,A\}}\ \bigcup_{\jmath\in\{0,A\}}
\mathcal Q_{\iota,\jmath}
\right)
\le
\sum_{\iota\in\{0,A\}}\ \sum_{\jmath\in\{0,A\}}
\mu(\mathcal Q_{\iota,\jmath})
\le\frac4{1600}.
\]
% lean: CausalSmith.Stat.FinitepHomogeneityDensegamma.ledgerCoeffError_indep
% lean: CausalSmith.Stat.FinitepHomogeneityDensegamma.ledgerCoeffError_cross_second_moment_le
% lean: CausalSmith.Stat.FinitepHomogeneityDensegamma.ledgerCoeffError_cross_tail

\item
Combining the two union bounds yields
\[
\mu(\mathcal F_{\mathrm{bad}})
\le
\frac4{800}+\frac4{1600}
=
\frac3{400}.
\]
% lean: CausalSmith.Stat.FinitepHomogeneityDensegamma.uniform_affine_profile_bound

\item
Fix a sample in \(\mathcal F_{\mathrm{bad}}^c\). By the definitions of the bad events,
\[
\|\operatorname{pr}_{\mathcal S_P}\mathsf E_{b,\iota}\|_2
\le40\sqrt{\Lambda_{n,v}},
\qquad
|\langle\mathsf E_{1,\iota},\mathsf E_{2,\jmath}\rangle|
\le40\sqrt M\,\Lambda_{n,v}.
\]
For every \(c\in\mathbb R\), the affine mean belongs to \(\mathcal S_P\). Orthogonal projection and Cauchy--Schwarz therefore give
\[
\begin{aligned}
|\langle\theta_P(c),\mathsf E_{b,\iota}\rangle|
&=
|\langle\theta_P(c),
\operatorname{pr}_{\mathcal S_P}\mathsf E_{b,\iota}\rangle|\\
&\le
\|\theta_P(c)\|_2\,
\|\operatorname{pr}_{\mathcal S_P}\mathsf E_{b,\iota}\|_2\\
&\le
40\sqrt{\Lambda_{n,v}}\,\|\theta_P(c)\|_2.
\end{aligned}
\]
Now take any \(|c|\le1/2\). The mean-plus-error expansion is
\[
W(c)-\|\theta_P(c)\|_2^2
=
\langle\theta_P(c),\mathsf E_1(c)+\mathsf E_2(c)\rangle
+
\langle\mathsf E_1(c),\mathsf E_2(c)\rangle.
\]
For the linear term, the triangle inequality yields
\[
\begin{aligned}
|\langle\theta_P(c),\mathsf E_1(c)+\mathsf E_2(c)\rangle|
&\le
\sum_{b=1}^{2}
\left(
|\langle\theta_P(c),\mathsf E_{b,0}\rangle|
+
|c|\,|\langle\theta_P(c),\mathsf E_{b,A}\rangle|
\right)\\
&\le
(2+2|c|)\,40\sqrt{\Lambda_{n,v}}\,\|\theta_P(c)\|_2\\
&\le
120\sqrt{\Lambda_{n,v}}\,\|\theta_P(c)\|_2.
\end{aligned}
\]
For the quadratic term, expansion of the two affine errors gives
\[
\begin{aligned}
|\langle\mathsf E_1(c),\mathsf E_2(c)\rangle|
&\le
|\langle\mathsf E_{1,0},\mathsf E_{2,0}\rangle|
+
|c|\,|\langle\mathsf E_{1,0},\mathsf E_{2,A}\rangle|\\
&\quad+
|c|\,|\langle\mathsf E_{1,A},\mathsf E_{2,0}\rangle|
+
|c|^2|\langle\mathsf E_{1,A},\mathsf E_{2,A}\rangle|\\
&\le
(1+2|c|+|c|^2)\,40\sqrt M\,\Lambda_{n,v}\\
&\le
\frac94\,40\sqrt M\,\Lambda_{n,v}
=
90\sqrt M\,\Lambda_{n,v}.
\end{aligned}
\]
Here \(|c|^2\le1/4\). Since both ledger terms are nonnegative,
\[
\begin{aligned}
|W(c)-\|\theta_P(c)\|_2^2|
&\le
120\sqrt{\Lambda_{n,v}}\,\|\theta_P(c)\|_2
+
90\sqrt M\,\Lambda_{n,v}\\
&\le
128\left(
\sqrt{\Lambda_{n,v}}\,\|\theta_P(c)\|_2
+
\sqrt M\,\Lambda_{n,v}
\right).
\end{aligned}
\]
The sample was fixed throughout and \(c\) was arbitrary in the profiling interval. Thus
\[
\mathcal F_{\mathrm{bad}}^c\subseteq\mathcal G_{\mathrm{prof}}.
\]
% lean: CausalSmith.Stat.FinitepHomogeneityDensegamma.ledgerTheta_inner_error_le_of_projection
% lean: CausalSmith.Stat.FinitepHomogeneityDensegamma.affine_profile_error_bound

\item
Finally, measurability of the bad event and the inclusion just proved imply
\[
\mu(\mathcal G_{\mathrm{prof}})
\ge
\mu(\mathcal F_{\mathrm{bad}}^c)
=
1-\mu(\mathcal F_{\mathrm{bad}})
\ge
1-\frac3{400}
=
\frac{397}{400}
=
0.9925.
\]
Together with the measurable-event argument and the zero-ledger case, this proves all the assertions.
% lean: CausalSmith.Stat.FinitepHomogeneityDensegamma.uniform_affine_profile_bound
\end{enumerate}
\end{proof}
